\section{RAG 与 Agent 的关系：不是替代，而是内化}

模型、应用和基础设施演进也会改变技术概念的可见度。本节以 RAG 为例说明：当一项能力从产品卖点沉淀为默认组件，讨论热度可能下降，但它在系统中的职责反而更稳定；Agent 只是把检索与规划、工具、记忆放进更长的闭环。

\subsection{RAG ``降温'' 的真实含义}
判断一项技术是否过时，不能只看社交媒体关键词频率，而要看其功能是否被更高层系统吸收。这里的核心问题是检索仍否承担事实接地、候选缩减和来源追踪；只要这些需求存在，RAG 的机制就不会因为名称降温而消失。
在问答环节，讲者回应了 ``RAG 是否过时''：RAG 仍然关键，只是被更深地内嵌到产品栈中。用户不再显式讨论它，就像很少有人讨论数据库，但所有系统都在用数据库。

\begin{importantbox}{RAG 的地位变化：从产品功能到系统部件}
在 Agent 系统里，检索能力常与工具调用、记忆机制、任务规划融合，形成 ``检索-推理-执行'' 的联合闭环，而非单一检索模块。
\end{importantbox}

\subsection{分阶段排序：成本与准确率的联合优化}
讲者使用推荐系统经验类比 RAG/检索流程：先用便宜模型做粗排，再用更强模型做精排。这个分层策略本质是 ``成本优先缩小候选集，再对高价值候选投入高精度算力''。

\begin{knowledgebox}{两阶段检索-重排策略}
\begin{itemize}
    \item 第一阶段：关键词、向量、规则过滤，目标是召回率与成本可控；
    \item 第二阶段：LLM 重排与上下文推理，目标是答案准确率与解释质量。
\end{itemize}
\end{knowledgebox}

\subsection{幻觉问题的边界：常识与事实的分界线}
讲者讨论了 ``模型记忆事实'' 与 ``上下文提供事实'' 的边界模糊性。常识和事实知识并非总能清晰切开，这也是幻觉长期难解的重要原因。

\begin{warningbox}{企业场景必须把事实来源显式化}
在企业系统中，答案质量不能依赖隐性记忆。需要可追踪数据源、版本化检索索引、可复现实验与可解释引用链路，才能把幻觉风险降到可管理范围。
\end{warningbox}

\subsection{本章小结}
RAG 没有过时，而是成熟并内化。Agent 时代的关键不是 ``是否使用 RAG''，而是 ``如何把检索、推理、执行放入同一可治理管线''。

